\documentclass[10pt,a4paper]{amsart}
\usepackage[margin=19mm]{geometry}
\usepackage{amsmath,amssymb,mathtools}
\usepackage[scr=rsfs,cal=euler]{mathalfa}
\usepackage{xfrac}
\usepackage[unicode,hidelinks,bookmarks=false]{hyperref}
\newcommand{\ie}{i.e.\ }
\newcommand{\cf}{cf.\ }
\newcommand{\resp}{respectively\ }
\newcommand{\Subsection}{\subsection}
\providecommand{\nobreakdash}{\nobreak\hbox{-}}
\setlength{\emergencystretch}{2em}
\multlinegap0pt
\usepackage[shortlabels]{enumitem}
\usepackage{amssymb}
\usepackage{tikz}
\usepackage{mathtools}
\usepackage{float}
\newtheorem{assumption}{Assumption}
\newtheorem{theorem}{Theorem}
\newtheorem{lemma}{Lemma}
\newcommand{\R}{\mathbb{R}}
\newcommand{\Rvalue}{\mathcal{R}}
\newcommand{\clball}{\overline{B}}
\newcommand{\dist}{\mathrm{dist}}
\newcommand{\dk}{d^k}
\newcommand{\pmin}{p_{\min}}
\newcommand{\retr}{\mathrm{R}}
\newcommand{\taug}{\underline{\tau}}
\newcommand{\tauk}{\tau_k}
\newcommand{\wk}{w^k}
\newcommand{\xk}{x^k}
\newcommand{\xkp}{x^{k+1}}
\title{A Nonmonotone Descent Method for Optimization Problems Defined by Upper-$\mathcal{C}^2 $ Functions over Submanifolds}
\author{Christian Kanzow, Leo Lehmann}
\date{Source: arXiv:2605.26909v1 (CC BY 4.0)}
\begin{document}
\maketitle

\noindent\emph{Source and licence.} A Nonmonotone Descent Method for Optimization Problems Defined by Upper-$\mathcal{C}^2 $ Functions over Submanifolds. Version arXiv:2605.26909v1 (CC BY 4.0), distributed by arXiv under the Creative Commons Attribution 4.0 International licence, http://creativecommons.org/licenses/by/4.0/, as recorded in the arXiv licence metadata for this exact version and shown on the abstract page https://arxiv.org/abs/2605.26909v1. Full licence notice: \texttt{licenses/arxiv-2605.26909-CC-BY-4.0.txt}.\par\smallskip

\noindent\textit{Editorial scope.} Counted unit: the article's theorem thm:convergencerateproof (source0.tex lines 1484--1494). The article's complete original proof of it is delivered with it (source0.tex lines 1496--1591, one \texttt{proof} environment of 96 lines). No mathematics is written, repaired or re-derived: the statement and the proof are the article's own lines, in the article's own notation, and no retained line is substituted or re-flowed.  Retained uncounted as the source-local context the proof needs: lines 123--125; lines 415--421; lines 538--550; lines 832--834; lines 854--857; lines 905--907; lines 1423--1425; lines 1464--1482.  Nothing was omitted from any retained span, so the package carries no omission record.  No retained line contains a citation, so the wrapper carries no bibliography and no bibliography entry.  No retained line contains a figure environment, an image-inclusion command or drawing code, so no image is delivered and the package declares no asset dependency.  Editorial formatting only, in addition: an AMS \texttt{amsart} preamble that restates the article's own declarations and macros used by the retained lines, namely the environments \texttt{assumption}, \texttt{theorem}, \texttt{lemma} on separate counters, together with the standard packages those lines need.  The retained environments print in this document as Assumption 1, Theorem 1, Lemma 1 in order of appearance, whereas the article prints the same items with the numbering of its own full document.  The complete native source of the article as distributed in the arXiv e-print for this exact version is bundled unchanged as \texttt{sources/201452/source0.tex}.
\begin{equation}\label{eq:prob}
	\min_{x \in \mathcal{M}} \varphi(x),
\end{equation}

\begin{assumption}\label{descentas}
	Let $\mathcal{M}$ be a manifold with retraction $\retr$ and assume that for all $\bar x \in \mathcal{M}$ there exists a constant $\tilde \kappa > 0$, $r>0$ and an open neighborhood $V\subseteq \mathcal{M}$ of $\bar x$, such that the objective function $\varphi: \mathcal{M} \to \R$ from \eqref{eq:prob} satisfies
	\begin{equation}\label{eq:descentpropm}
		\varphi \big(\retr_x(\xi_x)\big) - \big(\varphi(x) + \langle w_\mathcal{M}, \xi_x \rangle\big) \leq \tilde \kappa \|\xi_x\|^2,
	\end{equation}
	for all $x \in V$, $\xi_x \in \clball_r(0_x) \subset T_x \mathcal{M}$ and all $w_\mathcal{M} \in \partial_\mathcal{M} \varphi(x)$. Further, assume that $\partial_\mathcal{M} \varphi$ is locally bounded.
\end{assumption}

\begin{assumption}\label{genas}
Assume:
\begin{itemize}
	\item[(a)] There exists a constant $a > 0$ such that 
		\begin{equation}\label{eq:wdbound1}
			\langle \wk_\mathcal{M}, \dk\rangle \leq - a \|\dk\|^2, \text{ for all } k \in \mathbb{N}.
		\end{equation}
	\item[(b)] There exists a constant $b >0$ such that
		\begin{equation}\label{eq:wdbound2}
			\|\wk_\mathcal{M}\| \leq b \|\dk\|, \text{ for all } k \in \mathbb{N}.
		\end{equation}
\end{itemize}
\end{assumption}

\begin{equation}\label{eq:normxi}
	e \|\xkp -\xk\| \leq \Xi_k,
\end{equation}

\begin{equation}\label{distbound}
	\dist \big( 0, \partial \Phi(\xk) \big) 
	\leq \frac{2 b}{\taug_C} \| \xkp - \xk\|,
\end{equation}

\begin{theorem}\label{thm:convergence}
Let Assumptions~\ref{descentas} and \ref{genas} hold, let $\{\xk\}_K$ be a subsequence converging to some accumulation point $x^*$, and suppose that $\Phi:= \varphi + \delta_\mathcal{M}$ satisfies the KL property at $x^*$. Then the entire sequence $\{\xk\}$ converges to $x^*$.
\end{theorem}

\begin{equation}\label{eq:sumXis}
	\sum_{j=k_0-1}^k \Xi_j = \sum_{j=k_0}^{l(k_0)} \Xi_{j-1} + \sum_{j=l(k_0)}^k \Xi_j \leq \sum_{j=k_0}^{l(k_0)} (3 \Xi_{j-1} + \Xi_j) + 2 \hat c \sum_{j = k_0}^{l(k_0)}\chi \big( \Rvalue_j - \varphi(x^*) \big).
\end{equation}

\begin{lemma}\label{Lem:Aragon}
Let $ \{ s_k \} \subseteq [0, \infty )$ be any monotonically decreasing sequence satisfying $ s_k \to 0 $ and 
\begin{equation*}
	s_k^\alpha \leq \beta \big( s_k - s_{k+1} \big) \quad
	\text{for all $ k $ sufficiently large},
\end{equation*}
where $ \alpha, \beta > 0 $ are suitable constants. Then the following
statements hold:
\begin{itemize}
	\item[(a)] If $ \alpha \in (0,1] $, the sequence $ \{ s_k \} $
	converges linearly to zero with rate $ 1 - \frac{1}{\beta} $.
	\item[(b)] If $ \alpha > 1 $, there exists a constant $ \eta > 0 $
	such that 
	\begin{equation*}
		s_k \leq \eta k^{- \frac{1}{\alpha - 1}} \quad
		\text{for all $ k $ sufficiently large}.
	\end{equation*}
\end{itemize}
\end{lemma}

\begin{theorem}\label{thm:convergencerateproof}
Let Assumptions~\ref{descentas} and \ref{genas} hold, and suppose that $\{ \xk \}$ converges on some subsequence $\{ \xk\}_K$ to a limit point $x^*$ such that $\Phi$ has the KL property at $x^*$. Then the entire sequence $\{ \xk \}$ converges to $x^*$. Further, if the corresponding desingularization function is given by $\chi(t) = ct^{\theta}$ (for some $c>0$ and $\theta \in (0,1)$), then the following statements hold:
\begin{enumerate}[label=(\alph*)]
	\item If $\theta \in [1/2, 1)$, then $\{\Rvalue_k\}$ converges R-linearly to $\varphi(x^*)$ and $\{\xk\}$ converges R-linearly to $x^*$. \label{itm1}
	\item If $\theta \in (0, 1/2)$, then there exist constants $\eta_1, \eta_2 > 0$ such that for all $k$ large enough it holds that\label{itm2}
	\begin{align}
		\Rvalue_k - \varphi(x^*) &\leq \eta_1 k^{-\frac{1}{1-2\theta}},\\
		\|\xk - x^*\| &\leq \eta_2 k^{-\frac{\theta}{1-2\theta}}.
	\end{align}
\end{enumerate}
\end{theorem}

\begin{proof}
Taking Theorem~\ref{thm:convergence} into account, we only need to verify statements \ref{itm1} and \ref{itm2}. As a first step, let us prove the results for the sequence $\{\Rvalue_k\}$. We first claim that for $\theta \in (0,1)$ and with
\begin{equation*}
	\omega := \left(\frac{e \taug_C}{2 c \theta b}\right)^2,
\end{equation*}
it holds that
\begin{equation}\label{eq:sigma-inequality}
	\varphi(\xk) - \varphi(x^*) \leq \Big(\frac{1}{\omega}\Big)^{\frac{1}{2(1-\theta)}} \big( \Rvalue_k - \Rvalue_{k+1} \big)^{\frac{1}{2(1-\theta)}}
\end{equation}
for all $k \in \mathbb{N}$ sufficiently large. In fact, if $\varphi(\xk) \leq \varphi(x^*)$ holds, then the left-hand side of \eqref{eq:sigma-inequality}
is nonpositive, hence, the claim holds trivially.
By previous results, we may assume that $k$ is large enough such 
that $\xk \in \clball_\vartheta (x^*)$ and $\varphi(x^*) < \varphi(\xk) < \varphi(x^*) + \eta$ hold.

As $\Phi$ satisfies the KL property at $x^*$ with $\chi(t) = ct^{\theta}$ and as $\Phi(x) = \varphi(x)$ for all $x \in \mathcal{M}$, we have
\begin{align*}
	1 & \leq \chi' \big( \varphi(\xk) - \varphi(x^*) \big)
	\dist \big( 0, \partial \Phi(\xk) \big)\\
	& = c\theta \big( \varphi(\xk) - \varphi(x^*) \big)^{\theta -1}
    \dist \big( 0, \partial \Phi (\xk) \big).
\end{align*}
By \eqref{distbound}, this yields
\begin{equation*}
	1 \leq c\theta \frac{2 b}{\taug_C}  \big( \varphi(\xk) - \varphi(x^*) \big)^{\theta -1} \| \xkp - x^k \|,
\end{equation*}
which gives the inequality
\begin{equation}\label{eq:proofconvrate1}
	\|\xkp - x^k \| \geq \frac{1}{c \theta \frac{2 b}{\taug_C}}
	\big( \varphi(\xk) - \varphi(x^*) \big)^{1-\theta}.
\end{equation}
Further, from \eqref{eq:normxi}, we have
\begin{equation}\label{eq:proofconvrate2}
	\Rvalue_{k+1} - \Rvalue_k \leq - e^2 \|\xkp - \xk \|^2 .
\end{equation}
Combination of \eqref{eq:proofconvrate1} and 
\eqref{eq:proofconvrate2} yields
\begin{align*}
	\Rvalue_{k+1} - \Rvalue_k &\leq - e^2 \|\xkp - x^k \|^2\\
	 & \leq - e^2 \frac{1}{c^2 \theta^2 \frac{4 b^2}{\taug_C^2}} 
	\big( \varphi(\xk) - \varphi(x^*) \big)^{2(1-\theta)}\\
	& = -\omega \big( \varphi(\xk) - \varphi(x^*) \big)^{2(1-\theta)} .
\end{align*}
Rearranging these terms shows that the claim \eqref{eq:sigma-inequality}
holds.

Next recall that, by the acceptance criterion for the stepsize
$ \tauk $, we always have $ \varphi(\xkp) \leq \Rvalue_k $.
Hence, it follows that
\begin{equation}\label{eq:Rvalue-inequality}
	\Rvalue_{k+1} = (1-p_{k+1}) \Rvalue_k + p_{k+1} \varphi(\xkp) \leq (1-\pmin) \Rvalue_k + \pmin \varphi(\xkp).
\end{equation}
Denote by $\{s_k\}$ the sequence defined by $s_k := \Rvalue_k - \varphi(x^*)
\geq 0 $. Then $ s_k \to 0 $ monotonically, and we obtain
\begin{align*}
	s_{k+1} &\leq (1-\pmin) s_k + \pmin(\varphi(\xkp) - \varphi(x^*))\\
			&\leq (1-\pmin)s_k + \pmin \Big(\frac{1}{\omega}\Big)^{\frac{1}{2(1-\theta)}} \big( s_{k+1} - s_{k+2} \big)^{\frac{1}{2(1-\theta)}} ,
\end{align*}
where the first inequality follows from \eqref{eq:Rvalue-inequality}
and the second one results from \eqref{eq:sigma-inequality}.
	By the monotonicity of $\{s_k\}$, this implies
\begin{equation*}
	s_{k+2} \leq (1- \pmin) s_k + \pmin \Big(\frac{1}{\omega}\Big)^{\frac{1}{2(1-\theta)}} \big( s_{k} - s_{k+2} \big)^{\frac{1}{2(1-\theta)}},
\end{equation*}
which, in turn, yields
\begin{align*}
s_k &\leq \frac{1}{\pmin} (s_k -s_{k+2}) + \Big(\frac{1}{\omega}\Big)^{\frac{1}{2(1-\theta)}} \big( s_k - s_{k+2} \big)^{\frac{1}{2(1-\theta)}}\\
	&\leq \left(\frac{1}{\pmin} + \left(\frac{1}{\omega}\right)^{\frac{1}{2(1-\theta)}}\right) (s_k - s_{k+2})^{\min\{1, \frac{1}{2(1-\theta)}\}}
\end{align*}
for all $ k $ sufficiently large.
As for all $a, b > 0$ it holds that $1/\min\{a, b\} = \max\{1/a, 1/b\}$, it follows that
\begin{equation*}
s_k^{\max\{1, 2(1-\theta)\}} \leq \gamma (s_k - s_{k+2}),
\end{equation*}
where 
\begin{equation*}
\gamma := \left(\frac{1}{\pmin} + \left(\frac{1}{\omega}\right)^{\frac{1}{2(1-\theta)}}\right)^{\max\{1, 2(1-\theta)\}} > 0
\end{equation*}
is a constant.

By considering even and odd indices separately, we are now in the setting of Lemma~\ref{Lem:Aragon} and immediately obtain the 
corresponding rate-of-convergence results for the sequence $\{\Rvalue_k\}$ as $\theta \in (0, 1/2)$ implies that $2(1-\theta) > 1$ and $\theta \in [1/2, 1)$ implies that $2(1-\theta) \in (0, 1]$.

Let us now verify the statements for the sequence $\{\xk\}$. In view of Theorem~\ref{thm:convergence}, the equations in the proof of that result remain valid if $k_0-1$ is replaced by some $k$ sufficiently large. Note that $\Xi_{j-1} \leq \sqrt{\Rvalue_{j-1} - \varphi(x^*)} = \sqrt{s_{j-1}}$. Taking the monotonicity of the function $\chi$ and the sequences $\{\Rvalue_k\}$ and thus $\{s_k\}$ into account, it follows from \eqref{eq:sumXis} in combination with \eqref{eq:normxi} and the definition of $\chi$ that, for $l>k$:
\begin{align*}
	\|\xk - x^l\| &\leq \sum_{j=k}^{l-1} \|x^{j+1}-x^j\| \leq \frac{1}{e} \sum_{j=k}^{l-1} \Xi_j  \\
	&\leq \frac{1}{e} \sum_{j=k+1}^{l(k+1)} (3\Xi_{j-1} + \Xi_j) + \frac{2 \hat c}{e} \sum_{j = k+1}^{l(k+1)}\chi \big( \Rvalue_j - \varphi(x^*) \big)\\
	&\leq \frac{4}{e} \sum_{j=k+1}^{l(k+1)} \sqrt{s_{j-1}} + \frac{2 \hat c}{e} \sum_{j = k+1}^{l(k+1)}\chi ( s_j )\\
	&\leq \frac{4m}{e} \sqrt{s_k} + \frac{2 \hat c m}{e}\chi ( s_k )\\
	&\leq \tilde \eta s_k^{\min\{1/2, \theta \}}
\end{align*}
for all $ k $ sufficiently large, where 
\begin{equation*}
\tilde \eta := \frac{4 + 2 \hat c c}{e}m.
\end{equation*}
Taking now $l \to \infty$, together with the corresponding rate-of-convergence results for $\{s_k\}$ from the first part, this completes the proof.
\end{proof}

\end{document}
